% from ega1-5-fr.tex
\oldpage[I]{133}
est commutatif, et identifie $X\times_S Y$ au produit des
$(S\times_T S)$-préschémas $S$ et $X\times_T Y$, les projections
s'identifiant à $\pi$ et $(p,q)_T$.
\end{proposition}

En vertu de (\hyperref[I.3.4.3-fr]{3.4.3}), on est ramené à démontrer la
proposition correspondante \emph{dans la catégorie des ensembles}, en
remplaçant $X$, $Y$, $S$ par $X(Z)_T$, $Y(Z)_T$, $S(Z)_T$, $Z$ étant un
$T$-préschéma arbitraire. Mais pour la catégorie des ensembles, la vérification
est immédiate et laissée au lecteur.

\begin{corollary}[5.3.6]
\label{I.5.3.6-fr}
Le morphisme $(p,q)_T$ s'identifie (en posant $P=S\times_T S$) à
$1_{X\times_T Y}\times_P\Delta_S$.
\end{corollary}

Cela résulte de (\hyperref[I.5.3.5-fr]{5.3.5}) et
(\hyperref[I.3.3.4-fr]{3.3.4}).

\begin{corollary}[5.3.7]
\label{I.5.3.7-fr}
Si $f:X\to Y$ est un $S$-morphisme, le diagramme
\[
  \xymatrix{
    X\ar[r]^{(1_X,f)_S}\ar[d]_f &
    X\times_S Y\ar[d]^{f\times_S 1_Y}\\
    Y\ar[r]_{\Delta_Y} &
    Y\times_S Y
  }
\]
est commutatif, et identifie $X$ au produit des
$(Y\times_S Y)$-préschémas $Y$ et $X\times_S Y$.
\end{corollary}

Il suffit d'appliquer (\hyperref[I.5.3.5-fr]{5.3.5}) en remplaçant $S$ par
$Y$ et $T$ par $S$, et remarquant que $X\times_Y Y=X$
(\hyperref[I.3.3.3-fr]{3.3.3}).

\begin{proposition}[5.3.8]
\label{I.5.3.8-fr}
Pour que $f:X\to Y$ soit un monomorphisme de préschémas, il faut et il suffit
que $\Delta_{X|Y}$ soit un isomorphisme de $X$ sur $X\times_Y X$.
\end{proposition}

En effet, dire que $f$ est un monomorphisme signifie que pour tout
$Y$-préschéma $Z$, l'application correspondante
$f':X(Z)_Y\to Y(Z)_Y$ est une injection, et comme $Y(Z)_Y$ est réduit à un
élément, cela signifie qu'il en est de même de $X(Z)_Y$. Mais cela s'exprime
aussi en disant que $X(Z)_Y\times X(Z)_Y$ est canoniquement isomorphe à
$X(Z)_Y$, et le premier de ces ensembles étant $(X\times_Y X)(Z)_Y$
(\hyperref[I.3.4.3.1-fr]{3.4.3.1}), cela signifie que $\Delta_{X|Y}$ est un
isomorphisme.

\begin{proposition}[5.3.9]
\label{I.5.3.9-fr}
Le morphisme diagonal $\Delta_X$ est une immersion de $X$ dans
$X\times_S X$.
\end{proposition}

En effet, comme les applications continues $p_1$ et $\Delta_X$ des espaces
sous-jacents sont telles que $p_1\circ\Delta_X$ soit l'identité, $\Delta_X$
est un homéomorphisme de $X$ sur $\Delta_X(X)$. De même, l'homomorphisme
composé
$\mathscr{O}_x\to\mathscr{O}_{\Delta_X(x)}\to\mathscr{O}_x$ des
homomorphismes correspondant à $p_1$ et à $\Delta_X$ étant l'identité,
l'homomorphisme correspondant à $\Delta_X$ est surjectif~; la proposition
résulte donc de (\hyperref[I.4.2.2-fr]{4.2.2}).

On dit que le sous-préschéma de $X\times_S X$ associé à l'immersion
$\Delta_X$ (\hyperref[I.4.2.1-fr]{4.2.1}) est \emph{la diagonale} de
$X\times_S X$.

\begin{corollary}[5.3.10]
\label{I.5.3.10-fr}
Sous les hypothèses de (\hyperref[I.5.3.5-fr]{5.3.5}), $(p,q)_T$ est une
immersion.
\end{corollary}

Cela résulte de (\hyperref[I.5.3.6-fr]{5.3.6}) et de
(\hyperref[I.4.3.1-fr]{4.3.1}).

On dit (sous les hypothèses de (\hyperref[I.5.3.5-fr]{5.3.5})) que $(p,q)_T$
est l'\emph{immersion canonique} de $X\times_S Y$ dans $X\times_T Y$.

\begin{corollary}[5.3.11]
\label{I.5.3.11-fr}
Soient $X$, $Y$ deux $S$-préschémas, $f:X\to Y$ un $S$-morphisme~; alors le
morphisme graphe $\Gamma_f=(1_X,f)_S$ de $f$
(\hyperref[I.3.3.14-fr]{3.3.14}) est une immersion de $X$ dans
$X\times_S Y$.
\end{corollary}

C'est le cas particulier du cor. (\hyperref[I.5.3.10-fr]{5.3.10}) où l'on
remplace $S$ par $Y$ et $T$ par $S$
(cf. (\hyperref[I.5.3.7-fr]{5.3.7})).

